\section{The Stationary Cell Problem and the Field Limit}
\label{sec:cell}

The comparison estimate controls how two constrained fields respond to a
small change of law. In a bulk synchronous coupling the common Brownian noises
cancel, leaving the imposed relative drift and the averaging action of
edge contacts. We retain the order of those contacts in a stationary bulk
limit, then use the average transport of relative heights to identify the
macroscopic equation. A positive source will show that its coefficient
is a positive scalar, and exponential tilting will identify the Gaussian
field law.

Throughout this section lattice vertices have unit spacing:
$V_n\subset\Lambda$, where $\Lambda$ is the triangular lattice, and
$x\in V_n$ represents the physical point $x/n$. The vertex density is
$v=2/\sqrt3$. All bulk sampling regions have a fixed positive
macroscopic buffer from fixed pins and reflecting sign walls. Trace
estimates up to walls are separate from these bulk estimates.

Here are the quantitative comparison bounds used in the construction.
For a stationary synchronous difference $z=h-\widetilde h$, a fixed
short time interval has finite transmission clusters $C$, whose sizes
have all fixed moments under volume sampling. Write $e_C^2$ for their
quadratic loss, denoted by $D_C$ in the conservative comparison of
Section~\ref{s:coupling}. On each fixed interior macroscopic set $Q$,
\begin{equation}\label{b:energy-input}
 \E\sum_{C:C\cap nQ\ne\varnothing}e_C^2\le C_Q.
\end{equation}
The version with a drift difference $n^{-2}b(x/n)$ has an additional
uniformly bounded error. The filled differences $U_n$ obey, for every
$1<p<2$,
\[
 \sup_n\E\|U_n\|_{W^{1,p}(Q)}^p<\infty,
 \qquad
 \E\|U_n-z\|_{L^1(Q)}\longrightarrow0,
\]
where the second norm uses the original interpolated difference.
The $p$-th moment follows from the hole estimate and expected energy
bound after restricting the weight and hole moments to their
deterministic high-probability bounds; polynomial exceptional-event
bounds remove the restriction. A bound only on
$\E\|U_n\|_{W^{1,p}}$ would not suffice for the uniform-integrability
argument below. At a volume-sampled vertex, the spatial increments of
$z$ on every fixed finite neighborhood, multiplied by $n$, and the
corresponding scaled step displacements are tight. These estimates
remain local when the ambient domain is arbitrarily larger than the
comparison box, allow centering in a free box, and include small smooth
exponential tilts.

These bounds are supplied by Theorem~\ref{c:comparison}: initial slack
comes from \ref{f:sparse} and \ref{f:conditioned-sparse}, short-step
sparsity from \ref{f:step} and \eqref{c:short-step}, dissipation and
compactness from \eqref{c:energy-bound}--\eqref{c:compactness}, using
\ref{f:hole} and \ref{f:averages}. They apply to the bounded-pin and
ordinary one-sided sign-string laws specified there. Exact-value
adaptive freezing is used only with its full transcript and the
annealed marginal interpretation of Lemma~\ref{c:freezing}; preservation
of a marginal supplies no estimate for an arbitrary fixed transcript.
Shrinking two-sided reflecting intervals are excluded. The finite-volume
force covers and insulation are Lemmas~\ref{c:force} and
\ref{c:insulation}; Lemma~\ref{b:force-transfer} below passes their
conditional inequalities to the limiting forward filtration. The
ordinary fixed-boundary extension is Lemma~\ref{c:exterior-zero}.

\subsection{The bulk environment}

We first retain the driving Brownian motions as part of the local limit.
This is necessary because the tangent operator records the order of
contacts, not just the initial gradient configuration.

\begin{proposition}[Stationary bulk law]\label{b:bulk}
Under the inputs specified above, the gradient process and its driving
Brownian increments, viewed from a uniformly sampled vertex in a compact
macroscopic region separated by a fixed positive macroscopic distance
from all fixed pins and reflecting sign walls, have a unique local limit. This limit is
independent of the bounded boundary data and of the small smooth tilts under
consideration. It is stationary in space and time and spatially ergodic.
Its local trajectories are continuous and have normal reflection. On a
fixed short step, every vertex belongs to a finite transmission cluster;
there are almost surely no simultaneous distinct-face contacts. A face
visited away from the endpoints of a contact episode has positive
reflection impulse in that episode.
\end{proposition}

\begin{proof}
First consider bounded continuous cylinder functions of finitely many
gradients and Brownian increments in a fixed time window. Tile an interior
region by disjoint boxes of side $L$. In each box run a zero-boundary
equilibrium bath using the original Brownian drivers at its free sites.
The simultaneous stationary joining exists by time averaging the finite
reflected semigroup. The reference baths have independent path laws: their
drivers are independent, and the unique invariant marginal of their
product dynamics is the product of their equilibrium measures.

Discard a boundary layer of relative width $\eta$ in each box. The centered
free comparison, uniform in the ambient domain size, makes the difference
of each fixed local gradient tend to zero in volume probability as
$L\to\infty$. Stationarity gives the same assertion at any fixed finite
list of times. The drivers are identical. Consequently cylinder averages
for the original bath differ in $L^1$ from those for the references by a
quantity tending to zero, followed by an error $O(\eta)$ from discarded
sites. Comparing two original systems through reference boxes of the same
size proves that their volume-sampled cylinder laws are Cauchy.

For a spatial observation box of side $R\gg L$, its reference-box averages
are independent and bounded. Their average has variance at most
$C(L/R)^2$. Choose $L\to\infty$ slowly with $L/R\to0$, and then let
$\eta\downarrow0$. This proves concentration of spatial cylinder averages.
The estimates allow $R$ itself to grow arbitrarily slowly relative to the
original domain, so the concentration passes to the infinite-volume local
limit. Cylinder approximation then proves spatial ergodicity of the full
gradient-and-driver law. Spatial stationarity follows by shifting the
uniform root, and time stationarity passes from the finite stationary laws.
The same comparison includes the small smooth tilts and proves their
independence in the unscaled limit.

We give the path and reflection details separately from this static
comparison. For a stationary finite-volume bath, Lemma~\ref{f:step}
and the forward and reversed
martingale decompositions express each coordinate increment as half the
sum of two variance-two Brownian increments. The usual Brownian modulus
bound therefore gives tightness of each coordinate path after subtracting
its initial value. Relative initial heights in a fixed neighborhood are
bounded by graph distance. This gives tightness of all local continuous
paths, jointly with their drivers. A bounded drift of order $n^{-2}$
contributes a vanishing term. Independence of future Brownian increments
from the forward past passes first for finite cylinder tests and then,
by a monotone-class argument, for the full forward sigma-field.

Fix one short step and a finite set of vertices. Outside an event of
arbitrarily small probability, its incident transmission clusters lie in
a fixed finite neighborhood, and all edges crossing their boundary have
positive slack throughout the step. This is the short-time isolation input.
There are only finitely many possible clusters under this restriction.
Include their markings in the joint extraction. On the event corresponding
to a particular finite cluster $C$, the trajectory coincides with the
normally reflected motion using only the internal edge constraints of
$C$: removed edges exert no force. The finite-dimensional Skorokhod map is
continuous in the initial point and the common continuous drive on this
isolated problem. Use the orthonormal constant vector
$|C|^{-1/2}\mathbf1_C$ and an orthonormal basis of its orthogonal
complement. The constant coordinate is unconstrained variance-two
Brownian motion; every edge normal lies in the orthogonal complement,
where the Brownian covariance remains $2I$ and the edge polytope is
compact. The arithmetic mean, in contrast, has variance parameter
$2/|C|$.
Thus the limiting trajectory has the same normal reflection. Letting the
cluster restriction increase identifies the dynamics everywhere locally.

The improving initial strict-slack bounds imply that each fixed initial
edge is strict almost surely in the limit. On an isolated finite cluster,
a union of intersections of two distinct facets is polar for normal
reflection. One can verify this from the Neumann Dirichlet form: a
logarithmic cutoff in the two transverse coordinates has energy tending
to zero. Starting at an interior point, such a polar set is not reached.
This applies before restricting to the isolation event, so conditioning on
that future event is unnecessary. The countable collection of finite
clusters and time intervals proves absence of simultaneous contacts.
The one-dimensional normal coordinate and the strong Markov property at
its first hit show that a visited face has positive regulator increment
on an interval containing that hit in its interior. Applying this assertion
to rational enclosing intervals makes it simultaneous for all episodes
used below. Stationarity and strict initial slack also exclude contact at
either fixed step endpoint.
\end{proof}

\subsection{The tangent update}

Fix a sufficiently short step length $T$. We choose its transmission
partition as a function of the unmarked environment. Take a tolerance
$s>0$ below the one supplied by short-step sparsity, and join an edge
when its minimum slack during the step is at most $s$. Choose $s$ off
the atoms of the distributions of these minima; only countably many
values are excluded. The resulting components are finite with the stated
moments. On every fixed finite neighborhood their membership is stable
under local path convergence, except on a null event. This convention
also applies to the union of the two approximating replicas' edges.
It avoids retaining an unspecified auxiliary partition in the cell law.
For an oriented edge facet write
its normal as $a_e$; the sign does not affect the matrix
\[
 P_e=I-\frac{a_ea_e^{\mathsf T}}{|a_e|^2}.
\]
This matrix averages the two incident coordinates and fixes all others.
It is an orthogonal projection, is doubly stochastic, and preserves constants.

\begin{lemma}[A stable finite product]\label{b:tangent}
On a finite isolated cluster, consider a normal-reflection trajectory with
interior endpoints, no simultaneous distinct-face contacts, and positive
impulse in each visited-face episode. For two compared trajectories whose
initial states are sufficiently close to the reference initial state and
which use the same driving path, sufficiently close to the reference drive,
their synchronous difference update is exactly a finite product
\[
 J=P_{e_k}\cdots P_{e_1}.
\]
Consequently, if the initial relative differences multiplied by $n$ are
tight, their limiting update is $w'=Jw$ modulo constants. An additional
relative drift bounded by $Cn^{-2}$ makes a vanishing error after this
multiplication by $n$.
\end{lemma}

\begin{proof}
For each facet $e$, its contact times form a compact subset of $[0,T]$.
Different nonempty contact sets are disjoint and thus have positive
distance. Since the number of facets is finite, there can be only finitely
many switches of contact label. Repetitions of the same label can be grouped.
Choose a finite partition with endpoints off all facets so that each
contact interval allows only its designated facet, and each remaining
interval has no contact. All undesignated inequalities on each such
interval have positive minimum slack.

Continuity of the finite Skorokhod map preserves these inactive margins
under small perturbations. Positive impulse supplies a strict running-minimum
witness in the normal coordinate. Therefore both sufficiently close
synchronous paths reach the designated face within that interval. The
one-dimensional reflection formula then removes their normal difference
exactly; their tangential difference stays unchanged. Their difference at
the end of the interval is $P_e$ times its input difference. Contact-free
intervals preserve the difference. Composition proves the exact identity.

The conclusion is stronger than an error $o(1)$, which would not survive
multiplication by $n$. On each fixed regular path neighborhood the identity
is exact. Normal monotonicity bounds a finite-dimensional perturbation
from a differential drift by its integrated Euclidean magnitude, at most
$C(C)Tn^{-2}$. Its product with $n$ tends to zero.

For the random application, first bound the cluster size and the scaled
initial oscillation. Next bound the positive inactive margins, temporal
separations, and impulse witnesses away from zero. Each restriction has
probability tending to one when its bound is relaxed. The deterministic
argument applies on the restricted event. Remove the restrictions last.
The union of the two replicas' near-contact clusters provides the needed
isolation. Subtracting a common spatial constant is harmless because all
free edge constraints are invariant under that translation.
\end{proof}

Write $\mathcal E$ for the two-sided stationary bulk environment, including
the gradient trajectories and drivers. On each step, Lemma~\ref{b:tangent}
defines $J$ separately on its finite clusters. We now describe the class of
relative height solutions for which the mean flux will be unique.
The comparison bounds control filled relative gradients in $L^p$, whereas
the correction at a volume-sampled site is controlled by tightness.
We therefore record the large-scale slope by affine growth in volume,
which survives these stationary corrections. A slope is called attained
only when a stationary solution with the bounds below exists.

\begin{definition}[Controlled tangent solution]\label{b:controlled}
A controlled solution of slope $q\in\R^2$ is a space-time stationary
joining of $\mathcal E$ with real relative heights $(w_x(k))$, defined
modulo a spatial additive constant at each $k\in\mathbb Z$, such that:
\begin{enumerate}
\item The relative heights form a spatial cocycle: after choosing $w_0=0$,
they obey $w_{x+y}-w_x=w_y\circ\tau_x$, where $\tau_x$ shifts the joined
environment and relative field by $x$. They update by
$w(k+1)=J_k w(k)$ modulo constants.
\item For every $\delta>0$ there are spatial constants $c_R$ for which
\[
 \frac1{|\Lambda\cap B_R|}
 \#\{x\in\Lambda\cap B_R:
       |w_x-q\cdot x-c_R|>\delta R\}
 \longrightarrow0
\]
in probability as $R\to\infty$.
\item If $C_x$ is the step cluster of $x$, $H_x=|C_x|$, and
\[
 d_x=\frac1{|C_x|}\sum_{y\in C_x}
          \bigl(w_y^2-(Jw)_y^2\bigr),
\]
then $\E d_0<\infty$. For every fixed $a\ge0$,
$\E[H_0^a|(Jw-w)_0|]<\infty$.
\end{enumerate}
The loss defining $d_x$ is nonnegative and invariant under addition of
a constant. A slope is called attained if such a solution exists.
\end{definition}

\subsection{Extraction and growth of the relative heights}

\begin{lemma}[Volume-sampled tangent solutions]\label{b:extraction}
Let $z$ be a stationary comparison on scale $n$ covered by
Theorem~\ref{c:comparison}, with relative drift $n^{-2}b(x/n)$.
Sample $x$ in a fixed compact macroscopic region separated by a positive
macroscopic distance from all fixed pins and reflecting sign walls.
Extract local limits of
$n(z_{x+i}-z_x)$, jointly with $x/n$ and the bulk environment. Conditional
on the macroscopic position, the limit is a stationary spatial cocycle,
has update $w'=Jw$ modulo constants, and has affine volume growth with a
possibly random slope $q$. Conditional on almost every value of $q$, it is
a controlled solution on the original environment law. If the expected
macroscopic difference converges to $m$, then
\[
 \E[q\mid x/n=y]=\nabla m(y)
\]
in the distributional sense in the interior.
\end{lemma}

\begin{proof}
The comparison estimates give tightness on every finite neighborhood and
every finite set of step times. A diagonal extraction preserves the exact
additivity of the differences. Shifting the root by a fixed lattice vector
changes its macroscopic position by $o(1)$ and changes its sampling law
only by a vanishing boundary error. The conditional local law is therefore
spatially stationary. Time stationarity passes directly. Its environmental
marginal is the same ergodic law at every macroscopic position by
Proposition~\ref{b:bulk}. Lemma~\ref{b:tangent} identifies the update.

For integrability, distribute $n^2 e_C^2$ equally among the vertices of
$C$. Its expectation at a volume-sampled site is bounded by
\eqref{b:energy-input}. Cluster displacement is bounded by a fixed
polynomial in cluster size times the square root of cluster loss.
For any $1<r<2$, Holder's inequality consequently bounds the $r$-th
moment of each scaled displacement multiplied by any fixed cluster power.
Lower semicontinuity and finite-cluster testing give the corresponding
integrability for the limit.

Include a marking for one fixed sparse-hole construction in the extraction.
Randomize the block-grid phase, and choose a rim representative by a
translation-equivariant rule. By \ref{f:hole}, the filled limiting
potential $\widehat w$ has stationary gradients in $L^p$ for some $p>1$.
Let $\mathcal I$ be the invariant sigma-field of this joint spatial law.
The spatial ergodic theorem makes the rescaled gradients converge weakly
to their conditional mean $q$ given $\mathcal I$. To see that the potential
has affine growth, subtract its mean on each box and rescale by $R$.
Poincare's inequality and the spatial $L^p$ bound give weak
$W^{1,p}$ compactness and strong $L^p$ compactness. Every compact limit
has distributional gradient $q$, and is therefore affine. This proves
the asserted volume growth for $\widehat w$.

The correction $w_x-\widehat w_x$ is a finite stationary random variable:
the hole has finite diameter and every finite collection of original
scaled differences is tight. No integrability of this correction is
required to transfer volume growth. For fixed $\delta>0$, stationarity
and bounded convergence show that the expected fraction of sites where
$|w_x-\widehat w_x|>\delta R$ tends to zero. The original cocycle thus
has the same slope. The identical argument for its stationary step
displacement shows that $q$ is time invariant.

For any Borel slope event $B$, the function
$\Prob(q\in B\mid\mathcal E)$ is spatially invariant. Ergodicity of
$\mathcal E$ makes it constant. Conditioning on $q$ therefore leaves
the environmental marginal unchanged. Stationarity and the controlled
bounds disintegrate for almost every slope value, proving the claimed
deterministic-slope solutions.

Finally, the filled gradients are uniformly integrable in volume sampling.
Discrete summation by parts against a smooth macroscopic test identifies
their expected weak limit with $\nabla m$. The expected filled and original
potentials have the same limit by their vanishing $L^1$ difference.
The conditional spatial mean of the extracted filled gradients is $q$,
so its expectation is precisely the stated derivative.
\end{proof}

We next prove uniqueness at each attained slope. Two ingredients are
needed: dissipation removes a sublinear stationary difference, and the
actual contact graph connects the lattice. The latter is a statement about
the environment, not about any particular solution joined to it.

\begin{lemma}[Forward recurrence of contact tests]\label{b:recurrence}
Fix an edge. Suppose that there are forward-past measurable events $G_j$
with $\Prob(G_j)\to1$. For each $j$, suppose a finite measurable cover of
$G_j$ assigns to each class a fixed finite path of genuine free facets
connecting the edge endpoints. Assume each assigned facet can be forced
during a bounded time interval with conditional probability at least
$p_j>0$, given the entire forward history at the start of that interval.
Then the future contact graph connects the edge endpoints almost surely.
If these assumptions hold for every lattice edge, that graph is connected.
\end{lemma}

\begin{proof}
Use a time grid whose spacing accommodates a force attempt and is a
multiple of $T$. Choose measurably one class at an overlap. Maintain a
separate visit counter for each class and cycle through its assigned
facets. On infinitely many visits to that class, each facet is attempted
infinitely often. Conditional Borel--Cantelli, applied to the forward
filtration, makes each such facet occur infinitely often: its conditional
success probability at an attempted step is at least $p_j$.

Stationarity and Poincare recurrence imply
\[
 \Prob(\text{infinitely many visits to }G_j)\ge\Prob(G_j).
\]
Some class is visited infinitely often on this event, and its contact
path connects the prescribed endpoints. Let $j\to\infty$, and then
intersect the resulting probability-one events over the countably many
edges. Neither temporal ergodicity nor conditioning on a temporal
invariant event is used.
\end{proof}

We next verify the forward-past premise in the infinite-volume law.
This requires more than independence of future drives: finite-volume
insulation marks may depend on distant initial data.

\begin{lemma}[Passage of force opportunities to the bulk]
\label{b:force-transfer}
The finite-cover force and insulation estimates of
Lemmas~\ref{c:force} and \ref{c:insulation}, with the improving
good-configuration probabilities in \ref{f:conditioned-sparse} and
\eqref{c:short-step}, imply the hypotheses of
Lemma~\ref{b:recurrence} for the bulk environment.
\end{lemma}
\begin{proof}
Fix a lattice edge and a quality index $j$. Use the finite cover in
Lemma~\ref{c:force}, choosing the first neighborhood at an overlap.
It has finitely many labels $1,\ldots,M_j$, each carrying a fixed path
of genuine free facets joining the edge endpoints. Artificial grounding
is excluded by the collar slack, and there are no genuine pins or walls
in this bulk test box. Let $G_{j,n}$ mark the configurations satisfying
the initial slack and insulation conditions, and let $L_{j,n}$ be their
cover label. The estimates give
$\Prob(G_{j,n}=0)\le\delta_j+o(1)$ with $\delta_j\downarrow0$.
The force duration and its positive probability $p_j$ may depend on $j$.

At each time of a grid accommodating the force duration, these marks
are functions of the full finite-volume initial configuration, before
future drives are sampled. For an assigned facet $e$ write $C_e$ for
the event of visiting that facet during the next force interval.
The force and insulation proof, conditional on the full initial data,
and independence of subsequent Brownian increments from the past give
\begin{equation}\label{b:finite-force-inequality}
 \E\bigl[\Psi\mathbf1_{\{G_{j,n}=1,L_{j,n}=\ell\}}
                    \mathbf1_{C_e}\bigr]
 \ge p_j\E\bigl[\Psi\mathbf1_{\{G_{j,n}=1,L_{j,n}=\ell\}}\bigr]
\end{equation}
for every bounded nonnegative function $\Psi$ of the forward past.
Here the initial bad-insulation event was discarded before forcing;
no conditioning on a future cluster or future success event is used.

Retain the finite-valued marks in the local extraction constructing
the bulk law, for every fixed quality, spatial translate, and time $kT$.
For attempts of a particular quality use a coarser subgrid whose spacing
accommodates its force duration.
A diagonal extraction is possible because these index sets are
countable. The resulting marked extension is stationary. Its unmarked
marginal is still the unique bulk environment of
Proposition~\ref{b:bulk}. The limiting marks need not be functions of
the limiting initial state, and we do not assert that they are.

Take $\Psi$ to be a bounded nonnegative continuous cylinder function
of the unmarked past. Visiting a specified facet in a fixed compact
time interval is closed in the uniform local path topology: the
maximum of its oriented edge gradient equals the facet bound. The
discrete mark event is clopen. Thus the left integrand in
\eqref{b:finite-force-inequality} is upper semicontinuous, and
Portmanteau gives
\[
 \E\bigl[\Psi\mathbf1_{\{G_j=1,L_j=\ell\}}\mathbf1_{C_e}\bigr]
 \ge \limsup_n
 \E\bigl[\Psi\mathbf1_{\{G_{j,n}=1,L_{j,n}=\ell\}}
                    \mathbf1_{C_e}\bigr]
 \ge p_j\E\bigl[\Psi\mathbf1_{\{G_j=1,L_j=\ell\}}\bigr].
\]
The last expectation passes exactly by continuity. Approximation of
finite-dimensional past laws by continuous functions, followed by a
monotone-class argument over past cylinders, extends the inequality
to every bounded nonnegative function of the full unmarked forward
past $\mathcal F_0$. Contact at a deterministic interval endpoint is
a null event by Proposition~\ref{b:bulk}, so the closed-event passage
does not create an endpoint-only contact for the later tangent argument.

Define
\[
 a_\ell=\Prob(G_j=1,L_j=\ell\mid\mathcal F_0),
 \qquad a=\sum_{\ell=1}^{M_j}a_\ell.
\]
Dropping the mark from the left side of the preceding inequality gives
$\Prob(C_e\mid\mathcal F_0)\ge p_j a_\ell$ for every facet of
the path with label $\ell$. Since $\E(1-a)\le\delta_j$,
$\Prob(a<1/2)\le2\delta_j$. On $\{a\ge1/2\}$ choose a
maximizing label, with a fixed tie rule. Every facet of its path then
has conditional contact probability at least $p_j/(2M_j)$.
These are the required forward-past measurable opportunity classes.
Choose conditional-expectation versions equivariantly under the
countable spatial and time shifts; their joint law is stationary.
Their probabilities approach one, so Lemma~\ref{b:recurrence} applies.
The auxiliary marks served only in the transfer; the resulting
contact-connectivity assertion concerns the original environment alone.
\end{proof}

\begin{lemma}[Zero-slope uniqueness]\label{b:zero-slope}
A controlled tangent solution of slope zero is spatially constant.
\end{lemma}

\begin{proof}
Let $a$ be a nonnegative compactly supported Lipschitz function, equal to
one on the unit ball. For $M=\eps R$, let
\[
 \Phi_M(s)=
 \begin{cases}s^2,&|s|\le M,\\2M|s|-M^2,&|s|>M,\end{cases}
 \qquad
 F_R(w)=\min_c\sum_x a(x/R)\Phi_M(w_x-c).
\]
Its Lipschitz constant in each input coordinate is bounded by
$2M a(x/R)$. Hence $F_R(Jw)-F_R(w)$ is integrable, although the two
terms need not separately have finite expectations. They have the same
distribution by time stationarity. Truncate each scalar minimum at level
$K$ and use
$|(s\wedge K)-(t\wedge K)|\le|s-t|$; dominated convergence gives
\begin{equation}\label{b:stationary-loss}
 \E[F_R(Jw)-F_R(w)]=0.
\end{equation}

Choose a minimizing center $c_R$. On a cluster $C$ set
$a_C=\min_{x\in C}a(x/R)$ and
\[
 D_C^\Phi=\sum_{x\in C}
    \{\Phi_M(w_x-c_R)-\Phi_M((Jw)_x-c_R)\}\ge0.
\]
Clusterwise Jensen and evaluation of the new minimum at $c_R$ give
\begin{equation}\label{b:huber-cutoff}
 F_R(Jw)-F_R(w)
 \le-\sum_C a_CD_C^\Phi
       +C\eps\sum_{C:C\cap\operatorname{supp}(a(\cdot/R))\ne\varnothing}
             \sum_{x\in C}H_x^k|(Jw-w)_x|.
\end{equation}
Indeed the cutoff oscillation on $C$ is at most
$C\operatorname{diam}(C)/R$, while $\Phi_M$ is $2\eps R$-Lipschitz.
The controlled displacement moments bound the expectation of the last
sum by $C R^2$; clusters reaching the support from outside are covered
by a further fixed power of $H_x$. Equations
\eqref{b:stationary-loss}--\eqref{b:huber-cutoff} therefore bound the
normalized expected convex loss by $C\eps$.

For fixed $\eps>0$, volume sublinearity puts all but a vanishing fraction
of the weighted sites within $o(R)$ of a common center. The Huber
minimizer is itself $o(R)$ from that center. To verify this, the derivative
from an exceptional fraction $\delta$ of sites is at most
$2\eps R\delta$ times total weight; the derivative from the remaining
sites has a fixed sign and larger magnitude if the center is displaced
by more than their oscillation plus $C\eps R\delta$. Thus such a
displacement cannot minimize the loss.

First bound cluster size. The entire input of a typical interior cluster
then lies in the quadratic region, with probability tending to one.
Stochastic averaging keeps its output there. On that cluster
$D_C^\Phi=\sum_C(w_x^2-(Jw)_x^2)$, independently of the center.
The stationary integrable quadratic loss allows the exceptional fraction
to be discarded by truncating that loss before taking $R\to\infty$.
The cluster-size restriction is then removed using its moments. It follows
that $\E d_0\le C\eps$. Send $\eps\downarrow0$.

Each factor of $J$ is an orthogonal averaging projection. Its quadratic
loss is nonnegative, and the losses telescope. Zero total loss therefore
makes every constituent projection fix its input. The relative heights
are constant in time, and each actual contact equates its incident values.
Lemma~\ref{b:recurrence} connects every pair of neighboring sites and
proves spatial constancy. This argument never conditions on the tangent
solution, which may depend on future drives.
\end{proof}

\subsection{Mean flux and the macroscopic equation}

For a controlled solution define the cluster-averaged flux at $x$ by
\begin{equation}\label{b:flux-definition}
 \mathcal J_x(w)=
 -\frac1{T|C_x|}\sum_{y\in C_x}(y-\overline C_x)(Jw-w)_y,
 \qquad
 \overline C_x=\frac1{|C_x|}\sum_{y\in C_x}y.
\end{equation}
Conservation makes this independent of the chosen cluster origin, and the
controlled moments make it integrable. Write $j(q)=\E\mathcal J_0(w)$
when the slope is $q$. The minus sign is the sign of down-gradient transport.

\begin{proposition}[The cell flux and weak equation]\label{b:cell}
Under the preceding inputs, $j(q)$ is well defined on the vector space of
attained slopes and is linear there. If any nonzero slope is attained,
every slope in $\R^2$ is attained and $j(q)=Aq$ for a finite scalar $A$.
For a stationary comparison with drift difference $n^{-2}b(x/n)$,
every subsequential expected macroscopic difference $m$ satisfies
\begin{equation}\label{b:weak-equation}
 \int j(\nabla m)\cdot\nabla f=\int b f,
 \qquad f\in C_c^\infty(U),
\end{equation}
in each limiting open region $U$ free of observations.
\end{proposition}

\begin{proof}
Join two solutions of the same slope relatively independently over the
full two-sided environment. The joining is invariant under space and time
shifts. Their difference has zero slope; its dissipation and controlled
displacements are bounded by constant multiples of the sums for the two
solutions. Lemma~\ref{b:zero-slope} makes their relative heights identical.
This proves uniqueness of flux, including conditional uniqueness given
the environment. Addition of joined solutions and scalar multiplication
preserve the controlled class. Hence the attained slopes form a vector
subspace and the flux is linear on it.

The bulk law is invariant under lattice rotations and reflections, since
it is the unique limit of the symmetric zero-boundary model. A nonzero
attained slope and its sixty-degree rotation span $\R^2$. The resulting
linear flux map commutes with these rotations and reflections. A real
two-dimensional map commuting with a nontrivial sixty-degree rotation
is a scalar matrix plus a scalar multiple of the quarter-turn matrix;
a reflection eliminates the latter. Thus $j(q)=Aq$. Its coefficient is
finite by the integrability of \eqref{b:flux-definition}.

We justify passage of fluxes to expectations. At a volume-sampled site let
$D_n=n^2e_C^2/|C|$ and $H_n=|C|$. A scaled step displacement, or the
scaled cluster first moment per site, obeys
\[
 |F_n|\le C H_n^a\sqrt{D_n}
\]
for a fixed $a$. For $1<r<2$,
\begin{equation}\label{b:flux-ui}
 \E|F_n|^r\le
 C(\E D_n)^{r/2}
  (\E H_n^{2ar/(2-r)})^{(2-r)/2}.
\end{equation}
The right side is bounded. Thus these fluxes are uniformly integrable;
a moment of $D_n$ larger than one is not needed. Exceptional localization
events are removed using the polynomial error bounds in the comparison
estimate before applying this calculation.

Test stationarity of $z$ over one step against $f(x/n)$. On each complete
free cluster, the constant Taylor term sees only the added mass
$Tn^{-2}\sum_{x\in C}b(x/n)$, because reflections conserve the sum.
The first-order Taylor term is the negative of the cluster first moment
in \eqref{b:flux-definition}, with the height differences multiplied by
$n$ and with an overall factor $Tn^{-2}$. The second-order remainder has
expected magnitude bounded by
\[
 C n^{-3}\E\sum_{x\in n\operatorname{supp}(f)+O(H_x)}
                 H_x^a|n(z_T(x)-z_0(x))|=O(n^{-1}).
\]
The same moment estimate handles clusters touching the edge of the
testing region. The contribution of the added drift to the cluster
first moment is also $o(1)$. Lemmas~\ref{b:tangent} and
\ref{b:extraction}, finite-cluster truncation, and
\eqref{b:flux-ui} identify the limiting transport term as the conditional
mean of $j(q)$. The source and flux Riemann sums both carry the density
$v$, which cancels. Linearity and Lemma~\ref{b:extraction} give
$\E[j(q)\mid y]=j(\nabla m(y))$, proving
\eqref{b:weak-equation}. If the attained subspace is zero, this argument
still applies and its flux is identically zero.
\end{proof}

\subsection{Positive response and the limiting mean}

The scalar has not yet been shown nonzero. A positive source in a domain
with fixed zero boundary data supplies this information and, at the same
time, fixes the boundary condition in the response equation.

\begin{proposition}[Dirichlet response]\label{b:response}
Assume the comparison and average estimates hold up to identical fixed
boundaries of lattice domains converging to a bounded smooth domain $D$.
Then every slope is attained and the scalar in Proposition~\ref{b:cell}
satisfies $0<A<\infty$. For sufficiently small smooth sources $b$, the
mean response to relative drift $n^{-2}b(x/n)$ converges to
\[
 A^{-1}(-\Delta_D)^{-1}b.
\]
For boundary values $+\lambda$ and $-\lambda$ on the two marked arcs,
the untilted mean converges to their bounded harmonic extension.
\end{proposition}

\begin{proof}
Begin with zero fixed boundary values, and compare the untilted law with
its tilt by a nonzero nonnegative smooth source. The filled coupled
difference, extended by zero across its common pinned boundary, is bounded
in $W^{1,p}(\R^2)$ for some $p>1$, by
Lemma~\ref{c:exterior-zero}. To recall its geometric content, put the
original difference equal to zero on the boundary and exterior, and
retain the coercive edges incident to the common pins. Whenever a hole
meets this connected zero set, a path to infinity exits through a zero
rim vertex; choose that vertex as the hole's anchor. The usual rim
energy estimate holds for this anchor, so the extension is still bounded
in $W^{1,p}$ and remains exactly zero outside the polygon. Its expected
value is therefore weakly compact in that
space and strongly compact locally in
$L^p$. The original and filled expectations have the same limit.

Uniform convergence of the polygonal boundaries to the smooth boundary
implies that this zero extension vanishes outside an arbitrarily small
neighborhood of $\overline D$. Its limiting extension consequently
vanishes outside $D$. The Sobolev trace theorem for a smooth domain gives
$m\in W^{1,p}_0(D)$. Thus no assertion about random slit covariance is
needed for this ordinary boundary identification.

If only slope zero were attained, \eqref{b:weak-equation} would give
$\int b f=0$ for every compact test, a contradiction. A nonzero slope
is therefore attained, so the cell flux is $Aq$ on all of $\R^2$.
The same identity excludes $A=0$.

For completeness, the zero-trace weak equation has a unique
$W^{1,p}_0(D)$ solution even before an energy-space conclusion is known.
If $u$ is the difference of two such solutions, take a smooth compact
test $g$ and solve $-\Delta\phi=g$ with zero boundary data on smooth $D$.
The solution is smooth up to the boundary. Approximate $\phi$ in
$W^{1,p'}_0(D)$ by compactly supported smooth functions in the weak
harmonic equation, and integrate by parts using the zero trace of $u$.
This gives $\int u g=0$. Hence $m=A^{-1}(-\Delta_D)^{-1}b$.
Monotonicity under the nonnegative tilt gives $m\ge0$. The Dirichlet
Green solution of a nonzero nonnegative source is positive, so $A<0$
is impossible. This proves $A>0$.

For arbitrary identical bounded boundary data, the same zero-extension
argument and weak identity identify every mean-response subsequence with
$A^{-1}(-\Delta_D)^{-1}b$. Compactness therefore gives full convergence.

For the untilted mean, compare the prescribed-boundary bath with the
zero-boundary bath on the same domain. The latter has mean zero by height
reflection symmetry. Proposition~\ref{b:cell}, with source zero, makes
every limiting prescribed mean harmonic. Monotone comparison with the
zero-boundary law shifted by $\pm\lambda$ bounds this mean between
$-\lambda$ and $\lambda$. Near a compact subarc, compare instead with
the zero-boundary law shifted by that arc's prescribed value. The
coupled difference has exact zero boundary on that subarc. A cutoff
supported away from the marks permits the same local zero-extension
argument, giving the stated trace. A bounded harmonic function in smooth
$D$ is uniquely determined by these two open-arc values; the two marked
points have zero harmonic measure. This identifies the mean and every
subsequence.
\end{proof}

\subsection{Gaussian fluctuations}

\begin{theorem}[Field limit]\label{b:gff}
Under the inputs of this section for the ordinary fixed-boundary domains,
the centered affinely interpolated field converges as a random distribution
to $\sigma$ times the Dirichlet Gaussian free field with covariance
$(-\Delta_D)^{-1}$, where
\[
 \sigma^2=(vA)^{-1}.
\]
The stiffness is $vA$, and the limiting mean is the harmonic extension
identified in Proposition~\ref{b:response}.
\end{theorem}

\begin{proof}
For real $f\in C_c^\infty(D)$ define
\[
 X_n(f)=\frac1{vn^2}\sum_{x\in V_n}
             f(x/n)(h_x-\E h_x),
 \qquad L_n(t)=\log\E e^{tX_n(f)}.
\]
By \ref{f:averages}, there is $a_f>0$ such that
$\sup_n\E e^{a_f|X_n(f)|}<\infty$. Tilting by $e^{tX_n(f)}$
adds drift $n^{-2}tf/v$ to the variance-two reflected diffusion.
Proposition~\ref{b:response} therefore gives, for sufficiently small $t$,
\[
 L_n'(t)=\E_{n,t}X_n(f)
 \longrightarrow \frac{t}{vA}\int_D f(-\Delta_D)^{-1}f.
\]
For $|t|\le a_f/2$, Jensen's inequality bounds the denominator of the
tilted expectation below by one, while
\[
 |L_n'(t)|\le\E[|X_n(f)|e^{a_f|X_n(f)|/2}]
             \le C_{a_f}\E e^{a_f|X_n(f)|}.
\]
Integration from zero and dominated convergence yield the Gaussian log
moment-generating function near zero. The exponential bound gives
tightness and uniform integrability of these moment-generating functions;
uniqueness identifies every subsequential scalar law. Apply the same
argument to every finite real linear combination of tests to obtain
Gaussian finite-dimensional laws with covariance
$(vA)^{-1}(-\Delta_D)^{-1}$.

Here is the required distributional tightness. The anchored and
equal-mass average bounds in \ref{f:averages} give bounded second moments
for the coarse coefficients of a compactly supported wavelet expansion.
A normalized two-dimensional wavelet at spatial scale $2^{-j}$ has
coefficient equal to $O(2^{-j})$ times a smooth average difference.
Thus its variance is $O(2^{-2j})$ while there are $O(2^{2j})$
coefficients at that scale. Summation with weights $2^{-2sj}$ gives
a uniform local $H^{-s}$ second moment for each $s>0$, down to the
lattice scale. At finer scales the deterministic physical Lipschitz
bound $O(n)$ for the interpolant gives total squared coefficients
$O(n^2 2^{-2j})$, whose weighted tail is bounded as well. Compact
Sobolev embedding into a slightly more negative order and a countable
interior exhaustion give tightness as distributions.

Finally the Riemann pairing and affine pairing have vanishing difference.
An interior nodal hat function has integral $(vn^2)^{-1}$ and, by the
symmetry of its triangular star, zero first moment about its vertex.
Taylor expansion of $f$ therefore gives
\[
 \int f h^{\mathrm{aff}}-\frac1{vn^2}\sum_xf(x/n)h_x
 =O\!\left(n^{-2}\frac1{n^2}
                 \sum_{x\text{ near }n\operatorname{supp}(f)}|h_x|\right).
\]
Deterministic Lipschitzness and bounded boundary values give
$|h_x|=O(n)$ in the bounded rescaled domain, so this error is $O(n^{-1})$.
The same is true after centering. The scalar limits and tightness thus
identify the asserted affinely interpolated distributional limit.
\end{proof}
